% Generated draft: requires mathematical and source review.
\documentclass[11pt]{article}
\usepackage[margin=25mm]{geometry}
\usepackage{fontspec}
\setmainfont{Latin Modern Roman}
\usepackage{amsmath,amssymb,mathtools,unicode-math}
\setmathfont{Latin Modern Math}
\usepackage{enumitem,needspace,xurl,hyperref}
\setlength{\parindent}{0pt}
\setlength{\parskip}{5pt}
\setlength{\emergencystretch}{3em}
\newcommand{\N}{\mathbb{N}}
\newcommand{\Z}{\mathbb{Z}}
\newcommand{\Q}{\mathbb{Q}}
\newcommand{\R}{\mathbb{R}}
\newcommand{\C}{\mathbb{C}}
\newcommand{\F}{\mathbb{F}}
\newcommand{\A}{\mathbb{A}}
\newcommand{\PP}{\mathbb{P}}
\newcommand{\E}{\mathbb{E}}
\newcommand{\eps}{\varepsilon}
\newcommand{\dd}{\,\mathrm{d}}
\begin{document}
\section*{UnsolvedMath ID: 7200052}
\textbf{Source problem code:} catalogue entry
\subsection{Definitions and mathematical statement}
For a finite set $P\subset\mathbb R^d$, convex position means that every point is a vertex of $\operatorname{conv}(P)$, equivalently $p\notin\operatorname{conv}(P\setminus\{p\})$ for each $p\in P$. A projective transformation restricted to an affine chart has the form
\[
T(x)=\frac{Ax+b}{c^Tx+\delta},\qquad
\begin{pmatrix}A&b\\c^T&\delta\end{pmatrix}\in GL_{d+1}(\mathbb R),
\]
and is defined at each point of $P$ when $c^Tx+\delta\ne0$ there. The basic convexification condition is that such a $T$ exists with $T(P)$ in convex position.

The source names ``the McMullen problem on projectively transforming sets of points into convex position'' but gives no further statement. Determine the intended universal or extremal convexification assertion. A precise version must fix dimension and cardinality, any general-position hypothesis, whether all points or a selected subset must become vertices, and the bound or characterization being sought. The supplied summary speaks of subsets whereas the named item speaks of sets, so these versions are not silently identified. No particular extremal function or conjectured numerical value can be recovered from this record alone.
\subsection{Short English statement}
Study when a finite point set can be sent by a projective transformation into convex position. The source only names McMullen's problem and omits its dimension, cardinality, subset, and general-position quantifiers.
\subsection{Sources}
\begin{itemize}
\item Wikipedia contributors, \emph{List of unsolved problems in mathematics}, Geometry bullet 52, pinned revision 1366636767, \url{https://en.wikipedia.org/w/index.php?title=List_of_unsolved_problems_in_mathematics&oldid=1366636767#Geometry}; accessed 2026-07-29, MediaWiki source extraction. Wikipedia text is CC BY-SA 4.0, with attribution/share-alike handling required for release.
\item The supplied summary mentions Carath\'eodory/McMullen selection theorems without full references. No authoritative primary citation or exact McMullen theorem/conjecture statement is supplied.
\item Local registry \texttt{data/unsolvedmath/problems.json}, permanent UnsolvedMath ID $7200052$, source code catalogue entry, \emph{The McMullen problem on projectively transforming sets of points into convex position}. Dataset attribution: Ulam AI and the UnsolvedMath Contributors.
\end{itemize}
\end{document}
